\documentclass[a4paper,11pt]{article}

% =====================================================================
% 1. ENCODAGE, POLICES ET LANGUE
% =====================================================================
\usepackage[utf8]{inputenc}       % Gestion de l'UTF-8
\usepackage[T1]{fontenc}          % Encodage des caractères pour le français
\usepackage[french]{babel}        % Support du français
\usepackage{lmodern}              % Police propre et vectorielle
\usepackage{gensymb}

% =====================================================================
% 2. MISE EN PAGE ET GÉOMÉTRIE
% =====================================================================
\usepackage[a4paper, margin=3.2cm, top=3.5cm, bottom=4cm]{geometry}
\usepackage{setspace}             % Gestion de l'interligne
\usepackage{fancyhdr}             % En-têtes et pieds de page personnalisés (Remplace fancyheadings)
\usepackage{lastpage}             % Pour obtenir le numéro de la dernière page (\pageref{LastPage})
\usepackage[absolute,overlay]{textpos} % Positionnement absolu de blocs de texte
\usepackage{eso-pic}              % Pour ajouter des éléments en arrière-plan de la page

% =====================================================================
% 3. MATHÉMATIQUES ET SCIENCES
% =====================================================================
\usepackage{amsmath}              % Commandes mathématiques de base
\usepackage{amsfonts}             % Polices mathématiques
\usepackage{amssymb}              % Symboles mathématiques
\usepackage{siunitx}              % Alignement et écriture des unités scientifiques

% =====================================================================
% 4. GRAPHIQUES, TRACÉS ET DIAGRAMMES
% =====================================================================
\usepackage{graphicx}             % Inclusion d'images
\usepackage{svg}                  % Inclusion d'images SVG
\graphicspath{{Param/}, {Images/}} % Dossiers où chercher les images

\usepackage{tikz}                 % Dessin vectoriel personnalisé
\usetikzlibrary{shapes.geometric, arrows, chains, positioning, babel}
\usepackage{pgfplots}             % Tracé de graphiques mathématiques de haute qualité
\pgfplotsset{compat=1.18}

\usepackage[european, straightvoltages, RPvoltages]{circuitikz} % Tracé de circuits électriques

% =====================================================================
% 5. TABLEAUX
% =====================================================================
\usepackage{booktabs}             % Tableaux de qualité professionnelle (\toprule, \midrule, etc.)
\usepackage{tabularx}             % Tableaux à largeur fixe adaptative
\usepackage{array}                % Personnalisation avancée des colonnes
\usepackage{longtable}            % Tableaux multipages
\usepackage{makecell}             % Retours à la ligne et mise en forme dans les cellules
\usepackage{rotating}             % Pour tourner les tableaux ou les figures (mode paysage local)

% =====================================================================
% 6. STRUCTURATION, BLOCS ET COULEURS
% =====================================================================
\usepackage[table]{xcolor}        % Gestion des couleurs et coloration des lignes de tableaux
\usepackage{titlesec}             % Personnalisation des titres (\section, \subsection)
\usepackage{enumitem}             % Personnalisation des listes à puces
\usepackage{subcaption}           % Légendes pour les sous-figures
\usepackage{float}                % Forcer le positionnement des éléments graphiques ([H])
\usepackage{listings}             % Inclusion et coloration de code source (Python, C++, etc.)
\usepackage[most]{tcolorbox}      % Boîtes d'alerte et encadrés stylisés
\usepackage{pifont}               % Symboles et puces spécifiques (\ding)
\usepackage{lipsum}               % Générateur de faux texte de test

% =====================================================================
% 7. LIENS HYPERTEXTES (Toujours charger en dernier !)
% =====================================================================
\usepackage[hidelinks]{hyperref}  % Liens cliquables internes et externes sans bordures rouges

% Définition des couleurs
\definecolor{vhdlkeyword}{RGB}{0, 51, 153}      % Bleu foncé
\definecolor{vhdltype}{RGB}{153, 0, 153}        % Violet
\definecolor{vhdlstring}{RGB}{192, 0, 0}        % Rouge
\definecolor{vhdlcomment}{RGB}{34, 139, 34}     % Vert forêt
\definecolor{vhdlbackground}{RGB}{248, 248, 248} % Gris très clair

% Style VHDL sur mesure
\lstdefinelanguage{VHDL_Custom}{
  language=VHDL,
  morekeywords={
    architecture,entity,is,end,behavior,process,begin,
    signal,variable,constant,port,in,out,inout,buffer,
    generic,map,component,use,library,all,if,then,else,
    elsif,case,when,others,loop,while,for,generate,
    rising_edge,falling_edge,not,and,or,nand,nor,xor,xnor
  },
  morekeywords=[2]{
    std_logic,std_logic_vector,integer,boolean,natural,
    positive,character,string,signed,unsigned,bit,bit_vector
  },
  sensitive=false, % VHDL est insensible à la casse
  morecomment=[l]{--},
  morestring=[b]",
  morestring=[b]'
}

\lstdefinestyle{vhdlstyle}{
  language=VHDL_Custom,
  backgroundcolor=\color{vhdlbackground},
  basicstyle=\ttfamily\small,
  keywordstyle=\color{vhdlkeyword}\bfseries,
  keywordstyle=[2]\color{vhdltype}\bfseries,
  commentstyle=\color{vhdlcomment}\itshape,
  stringstyle=\color{vhdlstring},
  numbers=left,
  numberstyle=\tiny\color{gray},
  stepnumber=1,
  numbersep=8pt,
  showstringspaces=false,
  breaklines=true,
  frame=single,
  rulecolor=\color{black!20},
  tabsize=2,
  captionpos=b
}

\lstset{style=vhdlstyle}

\hypersetup{
    colorlinks=true,
    linkcolor=blue,
    citecolor=blue,
    urlcolor=blue,
    pdftitle={Rapport de projet},
    pdfauthor={Auteur},
    pdfkeywords={projet, rapport}
}

\begin{document}

\begin{titlepage}
    \centering

    % --- Logo en haut ---
    \includegraphics[width=0.9\linewidth]{Param/Polytech_Univ.png}\\[1.5cm]

    % --- Titre principal ---
    {\Huge \textbf{Rapport de Languages matériels et réutilisation}}\\[0.5cm]
    \rule{0.7\textwidth}{0.5pt}\\[0.8cm]
    {\Large \textbf{ETN5 - SETR}}\\[1cm]
    
    % --- Mise en page en deux colonnes (élève / encadrants) ---
    \noindent
    \begin{minipage}[t]{0.48\textwidth}
        \raggedright
        {\large \textbf{Élèves :}}\\[0.3cm]
        {\large
        Nolan \textsc{BUCHET}}\\[0.6cm]
        {\large
        Maxence \textsc{BOULEAU}}\\[0.6cm]
    \end{minipage}%
    \hfill
    \begin{minipage}[t]{0.48\textwidth}
        \raggedleft
        {\large \textbf{Date :}}\\[0.3cm]
        {\large
        Octobre 2026}
    \end{minipage}
    \vfill
\end{titlepage}

% LOGO IUT
\AddToHook{shipout/background}{%
  \begin{textblock*}{100cm}(100cm,100cm) % Ajustez les coordonnées pour décaler l'image
  \begin{tikzpicture}[remember picture, overlay]
    \node[anchor=north east, inner sep=18pt] at (current page.north east) {%
      \includegraphics[scale=0.18]{Logo_Universite_de_nantes.png}%
    };
  \end{tikzpicture}%
  \end{textblock*}
}


\AddToHook{shipout/background}{%
  \begin{textblock*}{100cm}(100cm,100cm) % Ajustez '15cm' pour déplacer l'image plus à droite
    \begin{tikzpicture}[remember picture, overlay]
      \node[anchor=north west, inner sep=25pt] at (current page.north west) {%
        \includegraphics[scale=0.50]{logo_poly.png}%
      };
    \end{tikzpicture}%
  \end{textblock*}
}


% NOM FICHIER
\AddToHook{shipout/background}{%
  \begin{textblock*}{1cm}(1cm,1cm)
  \end{textblock*}
  \begin{tikzpicture}[remember picture, overlay]
    \node[anchor=north, inner sep=50pt] at (current page.north) {%    
    {\fontsize{10}{13}\selectfont Rapport de Languages matériels et réutilisation}
    };
  \end{tikzpicture}%
}

% PAGE
\AddToHook{shipout/background}{%
  \begin{textblock*}{1cm}(1cm,1cm)
  \end{textblock*}
  \begin{tikzpicture}[remember picture, overlay]
    \node[anchor=south, inner sep=50pt] at (current page.south) {%
    \thepage/\pageref{LastPage}
    };
  \end{tikzpicture}%
}

% NOM PRENOM
\AddToHook{shipout/background}{%
  \begin{textblock*}{1cm}(1cm,1cm)
  \end{textblock*}
  \begin{tikzpicture}[remember picture, overlay]
    \node[anchor=south west, inner sep=50pt] at (current page.south west) {%    
    {\fontsize{10}{13}\selectfont Maxence BOULEAU \& Nolan BUCHET}
    };
  \end{tikzpicture}%
}

% DATE
\AddToHook{shipout/background}{%
  \begin{textblock*}{1cm}(1cm,1cm)
  \end{textblock*}
  \begin{tikzpicture}[remember picture, overlay]
    \node[anchor=south east, inner sep=50pt] at (current page.south east) {%    
    {\fontsize{10}{13}\selectfont Octobre 2026}
    };
  \end{tikzpicture}%
}


\setlength{\headheight}{14pt}
\pagestyle{fancy}

\fancyhead[L]{\large} % Left side of the header (section title)
\fancyhead[R]{\large} % Right side of the header (page number)

% Footer configuration
\fancyfoot[C]{\large} % Centered text in the footer

\renewcommand{\headrulewidth}{1pt} % Header rule thickness
\renewcommand{\footrulewidth}{1pt} % Footer rule thickness




\tableofcontents

\newpage

\section{Introduction}

Ce mini-projet vise à concevoir et valider par simulation un système numérique sur puce articulé autour 
d'un processeur RISC-V et d'une mémoire modulaire, tout en mettant en application les principes de généricité 
et de réutilisation de blocs matériels (IP cores). La démarche consiste d'abord à transformer une mémoire 
RAM élémentaire en VHDL pour la rendre paramétrable et à instancier un bloc mémoire étendu (quadram) associé 
à un décodeur d'adresses. Dans un second temps, le modèle de processeur décrit en Verilog est adapté pour 
supporter l'extension d'instructions RV32IM et pour s'interfacer avec la mémoire externe. L'ensemble est 
enfin réuni au sein d'une entité englobante mixte VHDL/Verilog puis validé sous ModelSim à l'aide de programmes 
de test en assembleur, illustrant ainsi l'intégralité du flot de co-conception et d'intégration système.
\\



\section{Récupération d’un modèle de RAM}

\subsection{Validation du fonctionnement de la RAM}

Nous récupérons le modèle \texttt{ramchip.vhd} et son testbench afin de pouvoir observer et analyser son comportement sur ModelSim.

% Injecter le graphe de fonctionnement de la RAM
\begin{figure}[H]
    \centering
    \includegraphics[width=1\textwidth]{1_RAM_modified.png}
    \caption{Chronogramme de la RAM}
    \label{fig:ram}
\end{figure}

Dans le chronogramme ci-dessus, nous observons que notre simulation représente correctement le comportement d'une RAM\@.
\\

Le bon fonctionnement se valide en activant \texttt{we}. À chaque front montant de \texttt{clock}, nous définissons l’adresse \texttt{address} de la donnée à écrire, puis nous plaçons cette donnée sur \texttt{datain}.
\\

\begin{itemize}
    \item À 20 ns : écriture de 01 à l’adresse 001 ; \texttt{dataout} restitue 01
    \item À 30 ns : écriture de 02 à l’adresse 002 ; \texttt{dataout} restitue 02
    \item À 40 ns : écriture de 03 à l’adresse 003 ; \texttt{dataout} restitue 03, et ainsi de suite
    \\
\end{itemize}

Un effet de décalage apparait, la valeur affichée correspond toujours à la donnée appliquée au cycle précédent.
Exemple entre $20\text{ ns}$ et $30\text{ ns}$, \texttt{address} vaut 001, \texttt{datain} vaut 01 et \texttt{dataout} restitue 00. Cet effet est normal car notre RAM est synchrone et ne restitue donc la donnée qu'au cycle suivant.
\\


\subsection{Modification du code de la RAM}

Dans un second temps, nous mettons à jour le code en remplaçant les \texttt{std\_logic\_vector} par des \texttt{unsigned} 
afin de pouvoir effectuer des opérations arithmétiques sur les adresses. Le fonctionnement reste identique à celui de la RAM initiale.
\\

\subsection{Ajout de la généricité de la RAM}

Pour rendre notre RAM plus générique, nous ajoutons des paramètres de taille de données et d'adresses.
Nous ajoutons donc les généricités \texttt{DATA\_LENGHT} et \texttt{ADDR\_LENGHT} qui permettent de définir la taille de la RAM\@.

\begin{lstlisting}[caption={Nouvelle entité de la RAM}, label={lst:ram_entity}]
entity sync_ram is
  generic (
    ADDR_LENGHT : integer := 10;
    DATA_LENGHT : integer := 8
  );
  port (
    clock   : in  std_logic;
    we      : in  std_logic;
    address : in  unsigned((ADDR_LENGHT - 1) downto 0);
    datain  : in  unsigned((DATA_LENGHT - 1) downto 0);
    dataout : out unsigned((DATA_LENGHT - 1) downto 0)
  );
end entity sync_ram;
\end{lstlisting}

\subsection{Ajout du signal de sélection}

L'ajout du signal \texttt{cs} permet de contrôler l'accès à la RAM et de placer sa sortie en haute impédance (\texttt{'Z'}) lorsqu'elle n'est pas sélectionnée, en vue de la création de la \texttt{quadram}.

\begin{lstlisting}[caption={Ajout du signal cs}, label={lst:ram_entity_cs}]
entity sync_ram is
  generic(DataLenght : INTEGER := 8; AddrLenght : INTEGER := 8);
  port (
    cs    : in std_logic;
    clock : in std_logic;
    ...
  );
end entity;

architecture RTL of sync_ram is
  ...
begin
  process(clock) is
  begin
    if rising_edge(clock) and cs = '1' then
      -- Ecriture synchronisee
    end if;
  end process;

  dataout <= ram(to_integer(unsigned(read_address))) when cs = '1' else (others => 'Z');
end architecture;
\end{lstlisting}

\begin{figure}[H]
    \centering
    \includegraphics[width=1\textwidth]{1_RAM_CS_modified.png}
    \caption{Chronogramme de la RAM avec le signal \texttt{cs}}
    \label{fig:ram_cs}
\end{figure}

La Figure~\ref{fig:ram_cs} présente le chronogramme de la RAM après l’ajout du signal de sélection \texttt{cs}. Lorsque \texttt{cs} vaut \texttt{'1'}, la RAM est sélectionnée : l’écriture est autorisée sur un front montant de \texttt{clock} si \texttt{we} est également actif, et la sortie \texttt{dataout} restitue la donnée mémorisée à l’adresse demandée. Lorsque \texttt{cs} vaut \texttt{'0'}, la mémoire n’est pas sélectionnée et sa sortie est placée en haute impédance (\texttt{'Z'}), afin de ne pas perturber un éventuel bus partagé. \\

Le chronogramme illustre le comportement synchrone de l’écriture : les données de \texttt{datain} sont enregistrées sur un front montant de \texttt{clock}, lorsque \texttt{cs} et \texttt{we} sont actifs. La sortie \texttt{dataout} est, quant à elle, conditionnée par \texttt{cs}. L’ajout de ce signal complète donc le fonctionnement de la RAM en permettant de contrôler son accès.

\end{document}